\documentclass[11pt, a4paper]{article}
\usepackage[a4paper, top=2.5cm, bottom=2.5cm, left=2cm, right=2cm]{geometry}
\usepackage{amsmath,amssymb,amsthm,microtype}
\usepackage{hyperref}

\theoremstyle{definition}
\newtheorem{problem}{Problem}
\theoremstyle{plain}
\newtheorem{theorem}{Theorem}
\newtheorem{lemma}{Lemma}

\begin{document}

\title{\vspace{-1.5cm}\textbf{Regional Quantitative Problem-Solving Circle} \\ \Large Problem Set 3: Number Theory \& Modular Congruence (Weeks 5--6)}
\author{\textbf{Lead Architect:} Mohamed Jad Srifi}
\date{\textbf{Term:} Fall 2025 -- Spring 2026}
\maketitle

\noindent\textbf{Instructions:} Do not rely on heuristic guessing. Every problem requires structural justification via Euclidean reductions, modular arithmetic properties, or Fermat's Little Theorem. Ensure all parity arguments are formalized.

\vspace{0.5cm}

\begin{problem}[Polynomial Divisibility \& Euclidean Reduction --- AIME Caliber]
Find all positive integers $n$ such that $(2n + 1) \mid (n^3 + 3n^2 + 7)$. 

\textit{(Hint: Utilize synthetic division or polynomial Euclidean reduction to isolate an integer remainder that is independent of the variable $n$.)}
\end{problem}

\vspace{0.5cm}

\begin{problem}[Order, Invariants \& Fermat's Little Theorem --- Olympiad Track]
Let $p$ be an odd prime. Prove that any prime divisor $q$ of the Mersenne number $2^p - 1$ must satisfy the congruence $q \equiv 1 \pmod{2p}$.

\textit{(Hint: Establish the multiplicative order of $2 \pmod{q}$, apply Fermat's Little Theorem, and leverage the parity of primes.)}
\end{problem}

\vspace{0.5cm}

\begin{problem}[Diophantine Integral Invariant --- MathMaroc / Putnam Style]
Determine all pairs of positive integers $(x, y)$ satisfying the Diophantine equation $x^2 - 3y^2 = 17$.

\textit{(Hint: Analyze the algebraic constraint under a small modulo, such as modulo 3, and evaluate the set of possible quadratic residues.)}
\end{problem}

\end{document}
